\section{Verification, scope and provenance}

The public entry point \nolinkurl{MainTheoremsPartII} exports
\nolinkurl{DaggerModelsPartIICompanion.mainA},
\nolinkurl{DaggerModelsPartIICompanion.mainB} and
\nolinkurl{DaggerModelsPartIICompanion.mainC}.
Their types are respectively the existing original Theorem A statement,
the full Theorem B statement with literal simplicial structure,
and the unchanged four-link Theorem C statement.
No theorem hypothesis, universe, model or adjoint was changed for this catalogue.

The catalogue has \CatalogueCount\ grouped entries and
\CatalogueDeclarationCount\ distinct advertised Lean declaration roots.
Its machine-readable source is
\nolinkurl{docs/formalization-companion/results.json}.
The selected-root audit is
\nolinkurl{verification/catalogue-20261125/CheckCatalogueAxioms.lean}.
Its saved result identifies the exact catalogue hash, source hashes, imported
artifact hashes, actual compiler exit and the printed transitive axiom lists.
The three main aliases were also freshly compiled together in the separate
repository layout.
Successful selected-root checks use only
\texttt{propext}, \texttt{Classical.choice} and \texttt{Quot.sound}.

Current original-source Lake content-rehash/no-build gates check the imported
project graph for the A/B proofs and the catalogue's native owner modules.
These gates verify source freshness against built imports without recompiling
the entire project.
The selected native Part II sources match the original workspace, and the
used Part I dependency source bytes retain their fixed revision.
The named C stages retain their original successful receipts and unchanged
source and axiom-log bytes.
Current separate-layout checks use verified existing import caches;
they are not cold source rebuilds.

The earlier A/B milestone recorded a successful 5,592-job aggregate and
standard-axiom audits of 3,400 Part II declarations and 2,675 B-owner declarations.
Those audit sets overlap and are not theorem counts.
The earlier snapshot differs from the current aggregate and two later frame
sources, so it is retained as historical evidence rather than silently reused
as certification of the entire current tree.
The current catalogue audit provides the separate present check of advertised
roots.

The original final C proof used 359 named stages with 1,182 printed standard
axiom reports.
Its 156 directly imported project roots passed a content-rehash/no-build gate.
The final C statement and its four literal component targets were also checked
separately and recompiled from saved sources.
These C-specific counts are not counts for the whole expanded catalogue.

Generic model-category factorization, lifting, localization, groupoid completion,
nerve and simplicial infrastructure are reused from Mathlib and the pinned
Part I development.
The Lean toolchain is \texttt{leanprover/lean4:v4.27.0};
the inherited Mathlib revision is
\nolinkurl{a3a10db0e9d66acbebf76c5e6a135066525ac900}.
Each entry distinguishes its construction and reused inputs.
The catalogue is not a claim of independent novelty for those foundations.

The Rezk recognition entry is a sufficient implication under zero-row component
surjectivity and all ordered-fiber weak equivalences.
It is narrower than the manuscript's full characterization by coherent-unitary
essential surjectivity.
The zero-vertex reflection entry retains its Quillen-adjunction premise.
No conditional strictification, decorated-evaluator or fixed-target-comma
interface is described as an unconditional proof of its missing input.
The original Theorem C is instead proved by the checked ordinary-completion
and exact raw-map route.
A general total-derived-functor API and completion of every older alternative
proof route are not claimed.

The source and axiom logs have been checked locally.
AI agents reviewed the main A/B/C statements and the selected result summaries;
that is not independent human certification.
The literal four-link audit is not an independent Mathlib-only Challenge.
No new hosted Linux CI, real Comparator replay with statement/definition
negative controls, or independent human correspondence review is claimed.
The workflow remains manual and is not triggered for this publication.

For a fixed paper reference, use the dedicated release tag
\texttt{\CompanionTag} or its exact commit.
The previous C-focused release remains a separate historical snapshot.
Release archives and their hash records bind the PDF, sources and verification
evidence to the published commit.
